\section{Literature Review}

\subsection{Prediction markets as information aggregation}

Prediction markets are commonly justified as mechanisms that aggregate dispersed information into continuously updated prices. The canonical review by \citet{wolfers2004prediction} documents applications in elections, macroeconomic forecasting, and organizational decision making and argues that market prices can provide useful forecasts. Long-run election-market evidence similarly demonstrates that such markets may be competitive with conventional forecasts \citep{berg2008accuracy}. These findings establish a serious informational case for prediction markets and caution against dismissing them as intrinsically uninformative.

The interpretation of a binary-contract price as a probability is nevertheless conditional. \citet{manski2006prediction} shows that equilibrium prices need not identify mean beliefs without restrictions on preferences, wealth, and belief distributions. \citet{wolfers2006interpreting} derive conditions under which prices may approximate mean beliefs, but their result does not imply that every displayed price is an objective probability. The distinction is central to the present study: an informative trading price can be publicly represented more broadly than its institutional and statistical foundations warrant.

\subsection{Collective belief, surveys, and representativeness}

A market price is also distinct from public opinion. It is formed by self-selected traders who differ in information, incentives, wealth, and risk preferences. \citet{dana2019markets} show that aggregated judgments elicited without trading can rival market forecasts in some settings, challenging the assumption that a traded price uniquely reveals collective belief. More generally, search behavior, social-media expression, headline tone, survey expectations, and market prices are different constructs. Treating them as interchangeable creates a construct-validity problem.

This project therefore does not evaluate ``public sentiment'' by correlating prices with Google search volume or generic social-media sentiment. Instead, it studies the public-representation claim itself: whether media and platforms describe a trader-generated price as the belief of a broader population.

\subsection{Manipulation and institutional conditions}

The accuracy and resilience of prediction markets depend on mechanism and participation. Experimental work finds that attempted manipulation may invite corrective trading under some conditions \citep{hanson2006information}, while other designs permit manipulators to influence decisions based on market prices \citep{deck2013affecting}. \citet{jian2012aggregation} further show that information aggregation varies with market structure. This literature rules out a simple binary conclusion that prediction markets are either manipulation-proof or inherently unreliable.

For the present design, manipulation and insider-information risks are treated as interpretive limitations that a public-facing text may disclose. Their presence does not automatically invalidate the cited price.

\subsection{Market quality}

Market microstructure distinguishes multiple dimensions of liquidity and price formation, including spread, depth, trading activity, price impact, and volatility \citep{ohara1995market,hasbrouck2007empirical,foucault2013market}. These dimensions matter because a visible price may be stale, costly to trade, or supported by little depth even when cumulative volume appears large. Prediction-market contracts introduce additional complications: bounded binary prices, changing time to resolution, heterogeneous resolution rules, and platform-specific price displays.

Accordingly, this study does not equate cumulative volume with quality. It measures market quality as a vector and treats any composite index as a robustness device rather than an observed ground truth.

\subsection{Framing, authority, and platform information}

Framing selects aspects of reality and makes them salient in ways that promote particular definitions, interpretations, and evaluations \citep{entman1993framing}. When a news report labels a contract price as a probability or public sentiment, it does more than repeat a number: it specifies what the number means and whose belief it represents. News institutions may also provide authority cues, although content analysis alone cannot establish that audiences accepted those cues.

The project is also connected to research on platform-generated information and social learning. Bondi's work on consumer ratings examines how apparently transparent aggregates can guide discovery while producing mislearning through selection and platform design \citep{bondi2019alone,bondi2023social}. The present study transfers this information-design problem to prediction-market prices circulated through media. It does not reproduce a review-rating model; it combines computational framing measurement with market microstructure.

\subsection{Research gap}

Existing research separately explains prediction-market forecasting, conditional probability interpretation, manipulation, market quality, media framing, and platform information design. What remains underdeveloped is a joint empirical account of how prediction-market prices are translated into public signals and whether the breadth of those claims is aligned with contract-level evidentiary quality.

The project addresses this gap by linking sentence-level representations to timestamp-aligned market snapshots. Its contribution is not the claim that prediction markets are inaccurate or that media partnerships necessarily corrupt journalism. It is a measurement framework for identifying configurations in which broad authority and representativeness claims coexist with low market quality and limited disclosure.
